%------------------------------------------------------------------------------%
\section{Outras Funções Trigonométricas}

Definimos as funções seno, cosseno e tangente que são as principais funções
trigonométricas, mas outras funções trigonométricas também são importantes.

\emph{Cosecante de $x$}
\[
  \csc x = \frac{1}{\sin x}, \quad\text{ se } \sin(x) \neq 0.
\]
\emph{Secante de $x$}
\[
  \sec x=\frac{1}{\cos x}, \quad\text{ se } \cos(x) \neq 0.
\]
\emph{Cotangente de $x$}
\[
  \cot x=\frac{1}{\tan x}, \quad\text{ se } \tan(x) \neq 0.
\]

\begin{example}
  Pesquise sobre o domínio, a imagem, o gráfico e a inversa de cada uma das
  funções definidas nessa seção.
\end{example}

\begin{example}
  Se \(\cos(x)= -\frac{2}{5}\) e $x$ no 3\textsuperscript{o} quadrante,
  calcule o valor de $\sec(x)+\tan(x)$.

  \solution
  Sempre que tivermos o valor de uma das funções trigonométricas $\sin(x)$ ou
  $\cos(x)$, a relação fundamental trigonométrica será útil para
  encontrar a outra função. Por exemplo, nesse exercício temos o valor de
  $\cos(x)$, então podemos encontrar o valor de $\sin(x)$ usando a
  relação fundamental trigonométrica. Assim,
  \begin{align*}
    \sin^2(x) + \cos^2(x)                     & = 1                       \\
    \sin^2(x) + \left( -\frac{2}{5} \right)^2 & = 1                       \\
    \sin^2(x) + \frac{4}{25}                  & = 1                       \\
    \sin^2(x)                                 & = 1- \frac{4}{25}         \\
    \sin^2(x)                                 & = \frac{21}{25}           \\
    \sin(x)                                   & = \pm \frac{\sqrt{21}}{5}
  \end{align*}
  Uma vez que $x$ no 3\textsuperscript{o} quadrante, então
  $\sin(x) = -\frac{\sqrt{21}}{5}$, assim:
  \[
    \sec(x)
    = \frac{1}{\cos(x)}
    =-\frac{5}{2}
  \]
  e
  \[
    \tan(x)
    = \frac{\sin(x)}{\cos(x)}
    =\frac{-\frac{\sqrt{21}}{5}}{-\frac{2}{5}}=\frac{\sqrt{21}}{2}
  \]
  Consequentemente
  \[
    \sec(x)+\tan(x)
    = -\frac{5}{2}+\frac{\sqrt{21}}{2}
    = \frac{-5+\sqrt{21}}{2}.
  \]
\end{example}

\begin{example}
  Se \(\sin(x) = \frac{3}{5}\) e $x$ no 1\textsuperscript{o} quadrante,
  encontre o valor da expressão
  \[
    z = \frac{\sec(x)- \csc(x)}{\tan(x)+ \cot(x)}
  \]
  Usando a relação fundamental trigonométrica temos:
  \begin{align*}
    \sin^2(x)+ \cos^2(x)                    & = 1               \\
    \left( \frac{3}{5}\right)^2 + \cos^2(x) & = 1               \\
    \cos^2(x)                               & = 1-\frac{9}{25}  \\
    \cos^2(x)                               & = \frac{16}{25}   \\
    \cos(x)                                 & = \pm \frac{4}{5}
  \end{align*}
  Sendo $x$ no 1\textsuperscript{o} quadrante, então
  \(\cos(x)= \frac{4}{5}\), assim:
  \begin{align*}
    \sec(x) & = \frac{1}{\cos(x)}       = \frac{5}{4};                                     \\[2mm]
    \tan(x) & = \frac{\sin(x)}{\cos(x)} = \frac{\sfrac{3}{5}}{\sfrac{4}{5}} = \frac{3}{4}; \\[2mm]
    \csc(x) & = \frac{1}{\sin(x)}       = \frac{5}{3};                                     \\[2mm]
    \cot(x) & = \frac{1}{\tan(x)}       = \frac{4}{3}.
  \end{align*}
  Portanto,
  \[
    z
    = \frac{\sec(x) - \csc(x)}{\tan(x) + \cot(x)}
    = \frac{\sfrac{5}{4} - \sfrac{5}{3}}{\sfrac{3}{4} + \sfrac{4}{3}}
    = -\frac{1}{5}.
  \]
\end{example}
